\begin{exsol}{18}{Gibbons Problem 2.1}{130}{Advanced}{L3}
\asked Prove Becker's Rotten Kid theorem: when a child's action $A$ produces incomes
$I_C(A)$, $I_P(A)$ and an altruistic parent then chooses a bequest $B$, the backwards-induction
outcome has the child choosing the $A$ that maximizes $I_C(A) + I_P(A)$.

\method Backward induction in a two-stage game with continuous actions (Gibbons \S2.1.A,
pp.~57--61; Tadelis \S8.3.2): solve the parent's problem for every pair of incomes, then the
child's problem given the parent's reply. The key step is that the parent's reply depends on the
incomes only through their sum, and raises the child's consumption when the sum rises.

\solution
\textbf{Notation and assumptions.} $A \ge 0$ the child's action; $I_C, I_P$ strictly concave with
maximizers $A_C, A_P > 0$; $B \in \mathbb{R}$ the bequest; $U, V$ increasing and strictly concave;
$k > 0$. Write $Y = I_C + I_P$ (family income) and $x = I_C + B$ (the child's consumption), so the
parent consumes $I_P - B = Y - x$. As is standard for this problem, $U$ and $V$ are
differentiable and the parent's problem has a solution.

\textbf{Step 1: the parent's problem depends only on $Y$.} Given $(I_C, I_P)$ the parent solves
\[
  \max_{B \in \mathbb{R}}\ V(I_P - B) + k\,U(I_C + B).
\]
Because $B$ is unrestricted in sign, choosing $B$ is the same as choosing $x = I_C + B \in \mathbb{R}$,
and the problem becomes
\begin{equation}
  \max_{x \in \mathbb{R}}\ W(x; Y) = V(Y - x) + k\,U(x). \label{eq:18W}
\end{equation}
$W$ is strictly concave in $x$ (a sum of strictly concave functions), so the maximizer $x^*(Y)$ is
unique and satisfies the first-order condition
\begin{equation}
  V'(Y - x^*) = k\,U'(x^*). \label{eq:18foc}
\end{equation}
The bequest is $B^*(I_C, I_P) = x^*(I_C + I_P) - I_C$: it depends on both incomes, but the child's
consumption depends only on their sum.

\textbf{Step 2: $x^*$ is strictly increasing in $Y$.} Take $Y' > Y$ and suppose
$x' = x^*(Y') \le x = x^*(Y)$. Then $Y' - x' > Y - x$. Strict concavity makes $V'$ strictly
decreasing and $U'$ decreasing, so
\[
  V'(Y' - x') < V'(Y - x) \qquad\text{and}\qquad k\,U'(x') \ge k\,U'(x).
\]
Applying \eqref{eq:18foc} at $Y$ and at $Y'$:
\[
  V'(Y' - x') = k\,U'(x') \;\ge\; k\,U'(x) = V'(Y - x) \;>\; V'(Y' - x'),
\]
a contradiction. Hence $x^*(Y') > x^*(Y)$. (With $U, V$ twice differentiable,
$dx^*/dY = V''/(V'' + kU'') \in (0,1)$: the parent passes part, not all, of an income gain to the
child.)

\textbf{Step 3: the child's problem.} Anticipating $B^*$, the child chooses $A \ge 0$ to maximize
\[
  U\bigl(I_C(A) + B^*(I_C(A), I_P(A))\bigr) = U\bigl(x^*(I_C(A) + I_P(A))\bigr) = U\bigl(x^*(Y(A))\bigr).
\]
$U \circ x^*$ is strictly increasing (Steps~1--2 and $U$ increasing), so $A$ maximizes the child's
payoff if and only if it maximizes $Y(A) = I_C(A) + I_P(A)$.

\textbf{Step 4: the maximizer exists and is unique.} $Y$ is strictly concave. For
$A < \min\{A_C, A_P\}$ both incomes are increasing, and for $A > \max\{A_C, A_P\}$ both are
decreasing, so $Y$ attains its maximum on the compact interval $[\min\{A_C,A_P\}, \max\{A_C,A_P\}]
\subset (0, \infty)$, at a unique point $A^*$. The backwards-induction outcome is
\[
  \bigl(A^*,\ B^*(I_C(A^*), I_P(A^*))\bigr), \qquad A^* = \arg\max_{A \ge 0}\ [\,I_C(A) + I_P(A)\,].
  \qquad\square
\]

\interp The child is selfish but faces a parent who ``taxes'' any loss of family income: if the
child raised $I_C$ at a larger cost to $I_P$, the parent, who cares about the split, would cut the
bequest by more than the child gained. Altruism on one side, plus the parent's last move, aligns the
first mover's incentives with joint surplus. The theorem needs the parent to move last with an
unrestricted transfer: if $B \ge 0$ were imposed and binding, or if the child moved after the
bequest, Step~1 would fail. Note that the child does \emph{not} pick $A_C$; it is between $A_C$ and
$A_P$.

\verify \texttt{check\_B\_L03.py} checks the sign of $dx^*/dY$ for several values of
$(V'', U'', k)$, and solves an example numerically: $U = V = -e^{-x}$, $k = 2$,
$I_C = 4A - A^2$ ($A_C = 2$), $I_P = 10A - A^2$ ($A_P = 5$). For each $A$ on a $0.05$ grid it finds
the parent's best bequest on a fine grid; the child's best action is $A = 3.5$, the maximizer of
$I_C + I_P = 14A - 2A^2$, not $A_C = 2$. The child's consumption matches the closed form
$x^* = (Y + \ln k)/2$, which is increasing in $Y$. Gibbons prints no answer.

\answer{The parent's best bequest gives the child consumption $x^*(I_C + I_P)$ that solves
$V'(Y - x) = kU'(x)$; it depends only on family income $Y$ and is strictly increasing in it. So the
child's payoff $U(x^*(Y(A)))$ is a strictly increasing function of $Y(A)$, and the child chooses
$A^* = \arg\max_A [I_C(A) + I_P(A)]$, which exists and is unique by strict concavity.}
\end{exsol}
